%==============================================================
% LAMPIRAN
%==============================================================

% Halaman standalone judul "LAMPIRAN" di tengah halaman
\clearpage
\vspace*{\fill}
\begin{center}
  {\fontsize{14pt}{17pt}\selectfont\bfseries LAMPIRAN}
\end{center}
\vspace*{\fill}
\addcontentsline{toc}{chapter}{LAMPIRAN}

%--------------------------------------------------------------
% LAMPIRAN 1: Konfigurasi Hyperparameter Model
%--------------------------------------------------------------
\clearpage

\refstepcounter{lampiran}
\addcontentsline{lamp}{lampiran}{\protect\numberline{\thelampiran}Lampiran 1\enspace
Konfigurasi \textit{Hyperparameter} Model TinyResNetV2}
\label{lamp:hyperparameter}

\noindent\hangindent=5.6em
\makebox[5.4em][l]{\textbf{Lampiran \thelampiran}}%
Konfigurasi \textit{Hyperparameter} Model TinyResNetV2 pada Eksperimen Terbaik

\vspace{4pt}

\centering
\begin{tabular}{l l}
  \toprule
  \textbf{Parameter} & \textbf{Nilai} \\
  \midrule
  Arsitektur          & TinyResNetV2 (3 stage residual) \\
  Lebar kanal         & 48 / 96 / 192 \\
  Ukuran input        & $224 \times 224$ piksel \\
  Jumlah kelas        & 5 \\
  Fungsi kerugian     & Focal Loss ($\gamma = 1{,}0$) \\
  Bobot kelas         & Berbanding terbalik dengan frekuensi \\
  Optimizer           & AdamW \\
  Laju belajar        & $10^{-3}$ \\
  Peluruhan bobot     & $5 \times 10^{-4}$ \\
  Penjadwalan LR      & \textit{Cosine annealing} \\
  Jumlah epoch        & 160 \\
  Ukuran batch        & 32 \\
  Augmentasi MixUp    & $\alpha = 0{,}2$, $p = 0{,}5$ \\
  Augmentasi CutMix   & $\alpha = 1{,}0$, $p = 0{,}5$ \\
  \textit{Label smoothing} & 0,1 \\
  \textit{Stochastic depth} & $0 \rightarrow 0{,}2$ \\
  Weight EMA decay    & 0,999 \\
  Protokol CV         & StratifiedGroupKFold, 5 \textit{fold} \\
  Platform komputasi  & Modal (A100-80GB) \\
  \bottomrule
\end{tabular}

\clearpage

%--------------------------------------------------------------
% LAMPIRAN 2: Konfigurasi Pipeline Segmentasi Pembuluh Darah
%--------------------------------------------------------------

\refstepcounter{lampiran}
\addcontentsline{lamp}{lampiran}{\protect\numberline{\thelampiran}Lampiran 2\enspace
Konfigurasi Pipeline Segmentasi Pembuluh Darah Terbaik}
\label{lamp:vessel_config}

\noindent\hangindent=5.6em
\makebox[5.4em][l]{\textbf{Lampiran \thelampiran}}%
Konfigurasi Pipeline Segmentasi Pembuluh Darah Terbaik (\textit{g40\_m60\_fine})

\vspace{4pt}

\noindent
\begin{tabular}{l l}
  \toprule
  \textbf{Parameter} & \textbf{Nilai} \\
  \midrule
  Kanal sumber & Kanal hijau (G) dari RGB \\
  CLAHE pertama & \textit{Clip limit} 6,0; \textit{tile} $16 \times 16$ \\
  Preprocessing top-hat & Kernel elips, jari-jari 5 piksel \\
  Gabor: sigma & 1,5; 2,5; 3,5; 5,0 piksel \\
  Gabor: lambda & 3; 5; 7; 10 piksel \\
  Gabor: sudut & 12 sudut ($0°$--$165°$, langkah $15°$) \\
  Median blur & Kernel $7 \times 7$ \\
  CLAHE kedua & \textit{Clip limit} 12,0; \textit{tile} $12 \times 12$ \\
  Meijering: sigma halus & 0,8; 1,4; 2,0; 2,8; 3,6; 4,5 piksel \\
  Formula fusi & $0{,}40 \times \text{gabor\_tophat} + 0{,}60 \times \text{meijering\_fine}$ \\
  Ambang batas & Persentil P0,16 dengan erosi FOV \\
  Seleksi komponen & 3 komponen terhubung terbesar \\
  Penutupan morfologis & Ellips $3 \times 3$, 1 iterasi \\
  \midrule
  \textbf{Hasil Dice (set pengujian)} & \textbf{0,4739} \\
  \bottomrule
\end{tabular}

\clearpage

%--------------------------------------------------------------
% LAMPIRAN 3: Distribusi Dataset Zhao2024
%--------------------------------------------------------------

\refstepcounter{lampiran}
\addcontentsline{lamp}{lampiran}{\protect\numberline{\thelampiran}Lampiran 3\enspace
Distribusi Lengkap Dataset Zhao2024}
\label{lamp:dataset}

\noindent\hangindent=5.6em
\makebox[5.4em][l]{\textbf{Lampiran \thelampiran}}%
Distribusi Lengkap Dataset Zhao2024 dan Statistik Kelompok

\vspace{4pt}

\noindent
\begin{tabular}{l c c c}
  \toprule
  \textbf{Kelas} & \textbf{Jumlah Citra} & \textbf{Persentase} & \textbf{Jumlah Kelompok} \\
  \midrule
  Laser Scars & 343 & 31,2\%  & 343 (singleton) \\
  Stage 3     & 261 & 23,7\%  & $\sim$120 \\
  Normal      & 236 & 21,5\%  & $\sim$100 \\
  Stage 2     & 165 & 15,0\%  & $\sim$80 \\
  Stage 1     & 94  &  8,6\%  & $\sim$79 \\
  \midrule
  \textbf{Total} & \textbf{1.099} & \textbf{100\%} & \textbf{$\sim$721} \\
  \bottomrule
\end{tabular}

\vspace{6pt}
\noindent Catatan: jumlah kelompok merupakan perkiraan berdasarkan analisis NCC $\geq$ 0,99.
Bekas luka laser (\textit{laser scars}) diperlakukan sebagai \textit{singleton group} karena
tidak memiliki kecocokan di kelas lain.
